\hypertarget{appendix-7a06e5baac04}{}
\section*{Copy-on-write}
\section{Copy-on-write}\label{copy-on-write}

\textbf{Curriculum depth:} Introduced · \textbf{First seen:} C003 · \textbf{Latest development:} C003

\subsection{Definition and mechanism}\label{definition-and-mechanism}

COW initially shares physical contents while deferring private copies. A permitted write to a COW-protected mapping faults; OS allocates/copies a frame, updates the writer's mapping, and retries. Other mappings retain old contents. Not every read-only page is COW, so not every write fault is recoverable.

\subsection{Prerequisites}\label{prerequisites}

\begin{itemize}
\tightlist
\item
  \hyperref[appendix-cc0909a7ef14]{Page faults and demand paging} --- Deferred private copying is implemented through a recoverable write fault.
\end{itemize}

\subsection{Used by}\label{used-by}

Update downstream uses when later chunks build on this concept.

\subsection{Recall and next depth}\label{recall-and-next-depth}

Show frame identity before and after one process writes.

\textbf{Pending:} Fork behavior, mapping flags, reference tracking and costs.

\subsection{Source and navigation}\label{source-and-navigation}

\href{Continuous\%20Notes/Computing\%20Foundations\%20-\%20From\%20Bits\%20to\%20C++\%20and\%20Systems\%20APIs\#Where\%20the\%20bytes\%20live}{C003 teaching chapter} · \hyperref[appendix-fd090d982a15]{Part definitions} · \href{}{Dependency Map (02-Foundations-II, appendix-4ad47bdd269a)}

\subsection{Reconstruction depth}\label{reconstruction-depth}

Copy-on-write delays physical copying until shared data is actually modified. At the virtual-memory level, two address spaces may initially map the same physical page as read-only. A write triggers a protection fault; the operating system allocates a private page, copies old contents, updates the writer's page table and resumes the instruction.

The key invariant is logical independence after a write, even though physical storage was shared before it. This is why process creation can be cheap initially without eagerly duplicating every writable page.

COW has costs that disappear in a high-level description: the first write can incur a page fault and page copy, dirty-page accounting changes, and large working sets can turn an initially cheap clone into substantial later memory traffic. Huge pages and NUMA placement make those costs workload-dependent.

Do not generalize OS page COW into a recommendation for C++ value types. Library-level copy-on-write strings or containers complicate iterator/reference validity, threading and mutation semantics. The mechanism is useful when ownership and fault or copy behavior are explicit.
